\documentclass[11pt,a4paper]{article}
\usepackage[margin=25mm]{geometry}
\usepackage{amsmath,amssymb,amsthm,hyperref,booktabs}
\setlength{\emergencystretch}{2em}
\newtheorem{lemma}{Lemma}
\newtheorem{theorem}[lemma]{Theorem}
\newcommand{\gen}[1]{\langle#1\rangle}
\newcommand{\diag}{\operatorname{diag}}
\newcommand{\sw}{\operatorname{Swap}}
\title{A new pairing for a core on eight points}
\author{Human Compression project: a component of the L044 recovery}
\date{}
\begin{document}\maketitle

Let the natural $A_{2m}$ have a string C-marking of rank $m$,
where $m\ge8$ is even. Let $H=G_i$ be a transitive maximal
omission preserving two blocks of size $m$, called its rows.
Every distinguished maximal omission $H_j=G_{i,j}$ is assumed
intransitive. Choose a mark $l$ interchanging the rows. Suppose
its two noncommuting neighbours in $H$ are $r,s$, in that order
around $l$, and let $C$ be the other $m-4$ marks. Thus $r,s$
commute, and every mark of $C$ commutes with $l$ and with at
least one of $r,s$.

We acquire the proved contracted-forest and local-support lemmas.
They supply one eligible selected edge for every mark other than
$l$, and the full graph $J$ of the contracted $r,s$ edges has
rank three. In this paper we impose the additional hypothesis
that $J$ is connected on exactly four $l$-pairs, denoted by $Y$.
All other $l$-pairs are isolated in $J$. The reduction to this
case is a separate obligation.

We assume the alternating rank bound $\lfloor(q-1)/2\rfloor$
in every degree $12\le q<2m$, and retain the exact acquired
smaller-degree maxima. These are the induction hypotheses of
the completed pair-block theorem. The complete transitive rank
theorem and extremal symmetric independent-generator theorem
are retained as well.

\begin{theorem}\label{main}
Under these hypotheses, if at least one of $r,s$ interchanges
the two rows, the configuration is impossible.
\end{theorem}

\section*{The selected bridges and the two actual forms}
Write the points as $x_0,x_1$, with $l$ interchanging the two
points in each pair. A mark in $C$ has the form
$\diag(\pi)l^\chi$, where $\pi$ is an involution of the pair
indices and $\chi\in\{0,1\}$. This is a genuine action for
$C$ marks; no such action is asserted for $r,s$.

\begin{lemma}\label{extras}
Each selected $C$ edge has an endpoint outside $Y$. Every other
contracted edge of that same mark lies in $Y$, and there is at
most one such additional edge.
\end{lemma}
\begin{proof}
By the acquired bridge lemma, $J$ together with the $m-4$
selected $C$ edges is connected, and every selected edge is a
bridge. No such edge can have both endpoints in the connected
graph on $Y$. Collapsing that graph to one vertex leaves a tree.

Take an additional contracted edge of a mark $g\in C$, and
join its endpoints by a simple path in the connected graph.
The path cannot contain a selected edge of another $C$ label:
the new edge would provide a bypass, which lifts using $l$ to
a walk in that label's actual omission, contradicting eligibility.
If the selected $g$ edge occurs, it cannot be the first or last
edge of this path, because the distinct transpositions of the
projected involution $g$ are disjoint. Both remaining path parts
then use only $J$, so both ends of that selected edge lie in
$Y$, already excluded. Hence the path uses only $J$, and the
additional edge lies in $Y$. The acquired local-support lemma
allows at most two simple edges or a triangle for one label.
A projected involution is a matching, so the triangle is impossible.
\end{proof}

Reverse the string if necessary so that $r$ interchanges rows.
For a permutation $p$ of the pair indices, define
\[
 \sw(p):x_0\longmapsto p(x)_1,\qquad
          x_1\longmapsto p^{-1}(x)_0.
\]
The acquired support lemma gives $r=\sw(p)$ for a three-cycle
$p$, with $p$ fixing all other indices. Name the four core pairs
$a,b,c,d$ so that $p=(a\ b\ c)$.

\begin{lemma}\label{forms}
After rotating $a,b,c$, precisely the following two forms remain:
\[
\begin{array}{ll}
 \mathrm M:&s=(b_0\ d_0)(c_1\ d_1),\\[2pt]
 \mathrm S:&s=\sw((b\ d\ c)).
\end{array}                                                   \tag{1}
\]
In both forms, the involution
\[
                       z=\diag((b\ c))l                       \tag{2}
\]
centralizes $r,l,s$, interchanges the rows, and equals $l$
outside $Y$.
\end{lemma}
\begin{proof}
If $s$ preserves the rows, it is a double transposition by the
acquired local-support lemma. Commutation with the row-swapping
$r$ forces one transposition in each row. If the first is $t$,
the second is $t^p$, using right permutation multiplication.
Since its edges and the triangle of $p$ have union $Y$, the
transposition $t$ joins $d$ to one of $a,b,c$. Rotate the latter
indices to obtain $t=(b\ d)$ and $t^p=(c\ d)$.

If $s$ interchanges the rows, the same support lemma gives
$s=\sw(q)$ for a three-cycle $q$. Its support meets that of
$p$ in two points, since their union has four points. Rotate
$a,b,c$ so that the support of $q$ is $\{b,c,d\}$. Commutation
is exactly $pq^{-1}=qp^{-1}$. Substituting the two orientations
of this three-cycle leaves $q=(b\ d\ c)$; the other orientation
does not satisfy the equation. Finally, direct substitution in
these two forms proves every assertion about (2).
\end{proof}

\section*{Marks that might break the new pairing}
Let $g=\diag(\pi)l^\chi$ be a mark in $C$. By
Lemma~\ref{extras}, write $\pi=e\tau$, where $e$ is its
selected transposition and $\tau$ is either the identity or a
single transposition in $Y$, disjoint from $e$.

Suppose first that $g$ commutes with $r$. The six points over
$a,b,c$ form the unique six-point orbit of $\gen{l,r}$; all
other orbits are pairs. Thus $e$ is either $(d\ x)$ or
$(x\ y)$, with $x,y$ outside $Y$. On $a,b,c$, the permutation
$\pi$ centralizes $p$ when $\chi=0$ and inverts $p$ when
$\chi=1$. Its involutory possibilities are therefore
\[
 (\tau,\chi)=(1,0),\quad ((a\ b),1),\quad
              ((a\ c),1),\quad ((b\ c),1).                  \tag{3}
\]
Only the middle two possibilities fail to commute with $z$.

If $g$ instead commutes with $s$ in form $\mathrm M$, its
selected edge is $(a\ x)$ or is entirely outside $Y$. Commutation
with the two row transpositions forces $(\tau,\chi)=(1,0)$
or $((b\ c),1)$: with no row interchange it preserves both
transpositions, while with row interchange it exchanges them.
Both possibilities commute with $z$.

In form $\mathrm S$, commuting with $s$ gives the analogue of
(3) on $b,c,d$. The two failures are $((b\ d),1)$ and
$((c\ d),1)$. Reverse $r,s$ and rename
\[
                  (a',b',c',d')=(d,c,b,a).
\]
This preserves the form $\mathrm S$ and changes these failures
into the middle cases of (3). Moreover conjugation by $z$ fixes
$r,l,s$ and exchanges $(a\ b)$ with $(a\ c)$, leaving either
possible selected edge unchanged. It therefore suffices in both
forms to exclude
\[
       g=\diag(e(a\ b))l,\qquad
       e=(d\ x)\ \hbox{or}\ e=(x\ y).                       \tag{4}
\]
Every mark in $C$ commutes with at least one of $r,s$, so this
list covers all possible failures of centrality with $z$.

\begin{lemma}\label{mixed}
In form $\mathrm M$, a mark (4) violates the actual intersection
property.
\end{lemma}
\begin{proof}
Put $s'=lsl$ and $t=lg$. The following uniform permutation
identities use right multiplication:
\[
 w=\begin{cases}
 (ts')^2(st)^2s(s't)^2ss's,& e=(d\ x),\\
 \bigl(l(sg)^2s\bigr)^2,& e=(x\ y).
 \end{cases}                                                \tag{5}
\]
In each case,
\[
                  w=(a_0\ b_0\ c_0)(a_1\ c_1\ b_1)=rl.
                                                               \tag{6}
\]
For completeness, substituting $s=(b_0\ d_0)(c_1\ d_1)$,
$s'=(b_1\ d_1)(c_0\ d_0)$ and the appropriate $t$
into (5) gives the point images
\[
\begin{array}{c|rrrrrr}
v&a_0&b_0&c_0&a_1&b_1&c_1\\\hline
v^w&b_0&c_0&a_0&c_1&a_1&b_1.
\end{array}
\]
The points over $d,x,y$ are fixed whenever present, and the
global row interchanges cancel on every other pair. This verifies
(6) on the entire domain. Hence $r=wl$ lies in $\gen{l,s,g}$,
contradicting IP for the singleton mark $r$ and the disjoint
marked set $\{l,s,g\}$. No finite enumeration is needed.
\end{proof}

\begin{lemma}\label{swapping}
In form $\mathrm S$, a mark (4) makes $H_l$ transitive.
\end{lemma}
\begin{proof}
The marks $r,s,g$ connect all eight core points by the walk
\[
 a_0\mathrel{\mathop{-}^{r}} b_1
 \mathrel{\mathop{-}^{s}}c_0
 \mathrel{\mathop{-}^{r}}a_1
 \mathrel{\mathop{-}^{g}}b_0
 \mathrel{\mathop{-}^{r}}c_1
 \mathrel{\mathop{-}^{s}}d_0
 \mathrel{\mathop{-}^{r}}d_1.
\]
All these are edges of the actual omission $H_l=\gen{r,s,C}$.
On every outside pair, $r$ equals $l$, so that pair also lies
inside an $H_l$-orbit. Consequently all $H_l$-orbits are unions
of whole $l$-pairs, and the eight core points lie in one orbit.
The connected graph consisting of $J$ and the selected $C$
bridges now joins all these orbits. Every edge of $J$ lies in
the core orbit, and each selected edge is an actual edge of a
mark of $H_l$. Thus $H_l$ is transitive, contrary to hypothesis.
\end{proof}

\section*{Maximality of the new pair system}
The preceding lemmas show that every mark of $H$ centralizes $z$.
Its orbits are pairs meeting the original rows once. In coordinates
matching those rows by $z$, the group $H$ therefore lies in
$\diag(S_m)\times\gen z$. Let $P$ be its transitive image
on the new pairs.

If all $m-1$ projected distinguished marks are independent, the
acquired extremal symmetric independent-generator theorem gives
$P=S_m$. Otherwise delete a redundant projected mark $j$.
The actual subgroup $H_j$ still projects onto $P$, but is
intransitive on the $2m$ points. Its two point orbits each meet
every new pair once. Hence the projection is faithful and gives
a genuine transitive string C-marking of $P$ of rank $m-2$.
For $m\ge10$, the complete acquired transitive rank theorem
forces $P=S_m$ or $A_m$, since $m-2>m/2+1$. The complete
degree-eight boundary excludes rank six for every proper
transitive group other than $A_8$. Thus $P$ is primitive in
all cases, proving that the new pair system is maximal.

The completed pair-block theorem for the parent $A_{2m}$ now
excludes the configuration under its explicit smaller-degree
induction. This proves Theorem~\ref{main}. The argument nowhere
requires $z$ to belong to $H$; centrality as a permutation and
the actual omission orbits suffice.

\paragraph{Attribution and scope.}
The historical target is part of Proposition~3.5 of
Cameron--Fernandes--Leemans--Mixer. The acquired forest and bridge
arguments were written by the root investigator, who also suggested
seeking a new central pairing here. This continuation supplies the
two actual forms, the uniform IP identities, the eight-point walk
and the maximality argument. The connected four-pair hypothesis
remains explicit; its reduction and the case where both neighbours
preserve rows are separate components. No whole-source recovery is
claimed by this paper.
\end{document}
